作为我国淡水渔业最发达、生物多样性最丰富的区域之一，长江流域长期承受过度捕捞、工业废水排放和水域塑料污染；葛洲坝、三峡大坝等水利工程又阻断了鱼类洄游通道。多种影响叠加后，渔业资源持续衰退，典型食物链结构发生退化。2021年1月1日，长江十年禁渔政策开始实施，此后局部水域水质改善，水草等水生植被逐步恢复。2025年监测数据显示，青、草、鲢、鳙四大家鱼资源量已恢复至禁渔前的约1.8倍，表明禁渔对基础资源和经济鱼类具有显著修复作用。
